% sections/09_conclusion.tex --- SkillStrike

We presented \sysname, a systematic study of \emph{skill composition attacks}---
a new class of agent supply-chain exploits in which an adversary
exploits the shared context of co-installed LLM skills to trigger
dangerous, undisclosed capabilities without user knowledge or consent.

Our three-stage framework establishes the attack from discovery through validation.
\lscps identifies over 4,000 exploitable source-sink pairs in \NskillsCH{} \clawhub skills
using \NpatternsLSCPS{}+ static patterns, with 100\% recall on
independently confirmed malicious skills versus 0\% for the platform's scanner.
LLM semantic verification (two-stage Codex chain-of-thought) improves precision
by $21\times$, reducing 4,154 LSCPS candidates to \NcodexConfirm{} confirmed
exploitable paths on \skillssh{}.
Live-agent experiments on two independent platforms confirm the attack:
on \clawhub, \NpairsCHshell/\NpairsCHshell{} shell-execution pairs and
\NpairsCHConeok/\NpairsCHCone{} LLM-context injection pairs succeed;
on \skillssh{}, \NpairsSSsuccess/\NpairsSSlaunched{} pairs across all five
sink capability families succeed~(\PairsSSrate),
with zero install-time or runtime safety gates triggered in either experiment.

Three structural conclusions follow.
First, the attack surface is \emph{compositional}:
it does not exist in any individual skill, only in combinations,
making per-skill auditing an insufficient defense.
Second, current agent safety policies are
heuristic-based and channel-dependent:
they block memory writes from direct search-context injection
but pass the same operations through alternate channels
and pass shell commands unconditionally through all channels.
Third, the attack is \emph{platform-agnostic}:
it replicates across two distinct agent runtimes, two marketplaces,
and two attack models (two-skill composition and self-exfil),
confirming it is a structural property of shared-context skill ecosystems,
not an artifact of any single platform's design.

LLM agent skill marketplaces are growing rapidly while operating without
capability isolation, permission manifests, or install-time composition audits.
As agent autonomy expands to financial operations, multi-agent orchestration,
and computer-use tasks, the consequences of this gap will compound.
We call on platform operators to adopt mandatory capability manifests,
runtime cross-skill context isolation, and publish-time static screening.
All tooling, datasets, and experiment artifacts are released to facilitate this work.

\paragraph{Acknowledgments.}
Omitted for anonymous review.
